% Options for packages loaded elsewhere
\PassOptionsToPackage{unicode}{hyperref}
\PassOptionsToPackage{hyphens}{url}
\PassOptionsToPackage{dvipsnames,svgnames,x11names}{xcolor}
%
\documentclass[
  11pt,
]{article}
\usepackage{amsmath,amssymb}
\usepackage{iftex}
\ifPDFTeX
  \usepackage[T1]{fontenc}
  \usepackage[utf8]{inputenc}
  \usepackage{textcomp} % provide euro and other symbols
\else % if luatex or xetex
  \usepackage{unicode-math} % this also loads fontspec
  \defaultfontfeatures{Scale=MatchLowercase}
  \defaultfontfeatures[\rmfamily]{Ligatures=TeX,Scale=1}
\fi
\usepackage{lmodern}
\ifPDFTeX\else
  % xetex/luatex font selection
\fi
% Use upquote if available, for straight quotes in verbatim environments
\IfFileExists{upquote.sty}{\usepackage{upquote}}{}
\IfFileExists{microtype.sty}{% use microtype if available
  \usepackage[]{microtype}
  \UseMicrotypeSet[protrusion]{basicmath} % disable protrusion for tt fonts
}{}
\makeatletter
\@ifundefined{KOMAClassName}{% if non-KOMA class
  \IfFileExists{parskip.sty}{%
    \usepackage{parskip}
  }{% else
    \setlength{\parindent}{0pt}
    \setlength{\parskip}{6pt plus 2pt minus 1pt}}
}{% if KOMA class
  \KOMAoptions{parskip=half}}
\makeatother
\usepackage{xcolor}
\usepackage[margin=1in]{geometry}
\usepackage{color}
\usepackage{fancyvrb}
\newcommand{\VerbBar}{|}
\newcommand{\VERB}{\Verb[commandchars=\\\{\}]}
\DefineVerbatimEnvironment{Highlighting}{Verbatim}{commandchars=\\\{\}}
% Add ',fontsize=\small' for more characters per line
\newenvironment{Shaded}{}{}
\newcommand{\AlertTok}[1]{\textcolor[rgb]{1.00,0.00,0.00}{\textbf{#1}}}
\newcommand{\AnnotationTok}[1]{\textcolor[rgb]{0.38,0.63,0.69}{\textbf{\textit{#1}}}}
\newcommand{\AttributeTok}[1]{\textcolor[rgb]{0.49,0.56,0.16}{#1}}
\newcommand{\BaseNTok}[1]{\textcolor[rgb]{0.25,0.63,0.44}{#1}}
\newcommand{\BuiltInTok}[1]{\textcolor[rgb]{0.00,0.50,0.00}{#1}}
\newcommand{\CharTok}[1]{\textcolor[rgb]{0.25,0.44,0.63}{#1}}
\newcommand{\CommentTok}[1]{\textcolor[rgb]{0.38,0.63,0.69}{\textit{#1}}}
\newcommand{\CommentVarTok}[1]{\textcolor[rgb]{0.38,0.63,0.69}{\textbf{\textit{#1}}}}
\newcommand{\ConstantTok}[1]{\textcolor[rgb]{0.53,0.00,0.00}{#1}}
\newcommand{\ControlFlowTok}[1]{\textcolor[rgb]{0.00,0.44,0.13}{\textbf{#1}}}
\newcommand{\DataTypeTok}[1]{\textcolor[rgb]{0.56,0.13,0.00}{#1}}
\newcommand{\DecValTok}[1]{\textcolor[rgb]{0.25,0.63,0.44}{#1}}
\newcommand{\DocumentationTok}[1]{\textcolor[rgb]{0.73,0.13,0.13}{\textit{#1}}}
\newcommand{\ErrorTok}[1]{\textcolor[rgb]{1.00,0.00,0.00}{\textbf{#1}}}
\newcommand{\ExtensionTok}[1]{#1}
\newcommand{\FloatTok}[1]{\textcolor[rgb]{0.25,0.63,0.44}{#1}}
\newcommand{\FunctionTok}[1]{\textcolor[rgb]{0.02,0.16,0.49}{#1}}
\newcommand{\ImportTok}[1]{\textcolor[rgb]{0.00,0.50,0.00}{\textbf{#1}}}
\newcommand{\InformationTok}[1]{\textcolor[rgb]{0.38,0.63,0.69}{\textbf{\textit{#1}}}}
\newcommand{\KeywordTok}[1]{\textcolor[rgb]{0.00,0.44,0.13}{\textbf{#1}}}
\newcommand{\NormalTok}[1]{#1}
\newcommand{\OperatorTok}[1]{\textcolor[rgb]{0.40,0.40,0.40}{#1}}
\newcommand{\OtherTok}[1]{\textcolor[rgb]{0.00,0.44,0.13}{#1}}
\newcommand{\PreprocessorTok}[1]{\textcolor[rgb]{0.74,0.48,0.00}{#1}}
\newcommand{\RegionMarkerTok}[1]{#1}
\newcommand{\SpecialCharTok}[1]{\textcolor[rgb]{0.25,0.44,0.63}{#1}}
\newcommand{\SpecialStringTok}[1]{\textcolor[rgb]{0.73,0.40,0.53}{#1}}
\newcommand{\StringTok}[1]{\textcolor[rgb]{0.25,0.44,0.63}{#1}}
\newcommand{\VariableTok}[1]{\textcolor[rgb]{0.10,0.09,0.49}{#1}}
\newcommand{\VerbatimStringTok}[1]{\textcolor[rgb]{0.25,0.44,0.63}{#1}}
\newcommand{\WarningTok}[1]{\textcolor[rgb]{0.38,0.63,0.69}{\textbf{\textit{#1}}}}
\usepackage{longtable,booktabs,array}
\usepackage{calc} % for calculating minipage widths
% Correct order of tables after \paragraph or \subparagraph
\usepackage{etoolbox}
\makeatletter
\patchcmd\longtable{\par}{\if@noskipsec\mbox{}\fi\par}{}{}
\makeatother
% Allow footnotes in longtable head/foot
\IfFileExists{footnotehyper.sty}{\usepackage{footnotehyper}}{\usepackage{footnote}}
\makesavenoteenv{longtable}
\setlength{\emergencystretch}{3em} % prevent overfull lines
\providecommand{\tightlist}{%
  \setlength{\itemsep}{0pt}\setlength{\parskip}{0pt}}
\setcounter{secnumdepth}{5}
\ifLuaTeX
  \usepackage{selnolig}  % disable illegal ligatures
\fi
\IfFileExists{bookmark.sty}{\usepackage{bookmark}}{\usepackage{hyperref}}
\IfFileExists{xurl.sty}{\usepackage{xurl}}{} % add URL line breaks if available
\urlstyle{same}
\hypersetup{
  pdftitle={EGCAR: exact computational reductions, memory use, and time-limited SGCA},
  colorlinks=true,
  linkcolor={Maroon},
  filecolor={Maroon},
  citecolor={Blue},
  urlcolor={Blue},
  pdfcreator={LaTeX via pandoc}}

\title{EGCAR: exact computational reductions, memory use, and
time-limited SGCA}
\author{}
\date{24 September 2026}

\begin{document}
\maketitle

\hypertarget{scope-and-what-was-actually-run}{%
\section{Scope and what was actually
run}\label{scope-and-what-was-actually-run}}

This note explains the computational implementation in \textbf{egcar
0.2.12}, and the SGCA elapsed-time change in \textbf{0.2.13}. It
distinguishes algebraic identities used in the code from proposed
improvements to comparison methods. The EGCAR entrywise and group
penalties remain separate estimators. The covariance divisor, loading
ridge, numerical floors, support thresholds, ADMM stopping rules, and
common cross-validation score were preserved.

\textbf{The earlier performance report used actual R runs of a dedicated
synthetic benchmark. It did not run the full all-methods example or the
full SLURM/CV sweep.} The script was
\texttt{inst/validation/benchmark\_0212\_worker.R}, launched by
\texttt{tools/benchmark\_0212.R}. Each timing is the median of three
runs. The ADMM iterations were capped deliberately at the same counts in
both versions, with zero stopping tolerances; these are fixed-work
comparisons, not time-to-convergence measurements. Package preparation
was excluded from the fit timer, while loading-factor construction was
included. The optional third-party comparison suite was skipped in that
validation.

The full mathematical specification remains the supplied
\texttt{gca\_scalability.pdf}. In that PDF, Algorithm 2 forms the
restricted covariance at \textbf{line 14} and requests eigenpairs at
\textbf{line 15}. Algorithm 1 requests eigenpairs at line 6.

\hypertarget{representation-and-preparation}{%
\section{Representation and
preparation}\label{representation-and-preparation}}

Let the centered views be \(X_k\in\mathbb R^{n\times p_k}\), let
\(p=\sum_kp_k\), and write

\[
\widehat\Sigma_{k\ell}=\frac{X_k^\top X_\ell}{n},\qquad
q=\sum_{k<\ell}p_kp_\ell.
\]

Only one coefficient block \(C_{k\ell}\in\mathbb R^{p_k\times p_\ell}\)
is stored for each unordered pair. The opposite block is represented by
its transpose and \(C_{kk}=0\). This avoids storing the complete
symmetric \(p\times p\) coefficient matrix. The edge blocks themselves
are dense arrays during fitting: coefficient sparsity does not
automatically imply a sparse numerical representation.

The fitted loss is

\[
\ell(C)=\sum_{k<\ell}\left\{
\frac12\operatorname{tr}(C_{k\ell}^\top\widehat\Sigma_{kk}
 C_{k\ell}\widehat\Sigma_{\ell\ell})
-\langle\widehat\Sigma_{k\ell},C_{k\ell}\rangle\right\}.
\]

The two objectives are \(\ell(C)+\rho\sum_{k<\ell}\|C_{k\ell}\|_1\) and
\(\ell(C)+\lambda_g\sum_k\|M_k(C)\|_{2,1}\), respectively, with the
paper's concatenation

\[
M_k(C)=[C_{k1},\ldots,C_{k,k-1},C_{k,k+1},\ldots,C_{kK}].
\]

The optimized implementation uses this concatenation mathematically but
does not assemble it during ordinary group iterations.

\hypertarget{reuse-covariance-factorizations}{%
\subsection{Reuse covariance
factorizations}\label{reuse-covariance-factorizations}}

For each view, preparation computes a spectral representation

\[
\widehat\Sigma_{kk}=Q_kD_kQ_k^\top,
\qquad D_k=\operatorname{diag}(d_{k1},\ldots,d_{ka_k}).
\]

When \(n<p_k\), the code obtains this from a thin SVD of \(X_k\), with
\(a_k\leq\min(n,p_k)\), rather than computing an otherwise unused full
covariance eigendecomposition. It retains the positive covariance range
under the existing numerical convention; it does \textbf{not} truncate
that range to the requested loading rank \(r\). In exact arithmetic,
centering further bounds the rank by \(n-1\).

The bases, pairwise products \(d_{ki}d_{\ell j}\), projected
cross-covariances, and multiplication-order decisions are prepared once.
They are reused over iterations and over a warm-start penalty path
within the same training fold. Fold-specific matrices are never shared
across different training splits as though they were identical.

The current prepared object still contains the dense within-view and
cross-view covariance blocks. Thin spectral preparation saves
eigensystem work and storage; it does not make all EGCAR storage linear
in \(p\).

\hypertarget{relation-to-ccar3}{%
\subsection{Relation to ccar3}\label{relation-to-ccar3}}

The ccar3 manual and \texttt{R/ecca.r} supplied useful computational
patterns: thin data SVDs, cached spectral denominators,
reduced-coordinate updates, choosing matrix multiplication order, and
warm starts along a penalty path. These patterns informed the
implementation audit. EGCAR's compressed group-consensus state and final
symmetric operator callback are derived below for its own objective. We
did not substitute ccar3's two-view estimator, spectral cutoff, or
full-decomposition fallback for EGCAR's statistical rules. See sources
S1 and S2.

\hypertarget{the-coefficient-solve-an-exact-reduced-calculation}{%
\section{The coefficient solve: an exact reduced
calculation}\label{the-coefficient-solve-an-exact-reduced-calculation}}

Both EGCAR families repeatedly solve an equation of the form

\[
\widehat\Sigma_{kk}C_{k\ell}\widehat\Sigma_{\ell\ell}
+\alpha C_{k\ell}
=\widehat\Sigma_{k\ell}+\alpha T_{k\ell}.
\tag{1}
\]

For L11, \(\alpha=\mu\) and \(T=Z-V\), where \(V\) denotes its scaled
dual variable. The source names that L11 dual array \texttt{H}. For L21,
\(\alpha=2\mu_g\) and

\[
T_{k\ell}=\frac12\left\{
E_{k\ell}^{(k)}(G_k-V_k)+E_{k\ell}^{(\ell)}(G_\ell-V_\ell)
\right\}.
\tag{2}
\]

The factor two must remain: each edge participates in two
group-consensus constraints.

Define

\[
\widetilde T=Q_k^\top T_{k\ell}Q_\ell,
\qquad \widetilde\Sigma=Q_k^\top\widehat\Sigma_{k\ell}Q_\ell,
\qquad R_{k\ell}=\widehat\Sigma_{k\ell}-Q_k\widetilde\Sigma Q_\ell^\top.
\]

Then form the small array

\[
\widetilde C_{ij}
=\frac{\widetilde\Sigma_{ij}+\alpha\widetilde T_{ij}}
       {d_{ki}d_{\ell j}+\alpha},
\tag{3}
\]

and recover the full edge by

\[
C_{k\ell}=T_{k\ell}+\frac{R_{k\ell}}{\alpha}
+Q_k(\widetilde C-\widetilde T)Q_\ell^\top.
\tag{4}
\]

Equation (4) is essential when the covariances are singular.
Multiplication by the covariance blocks acts only on the doubly
projected component. On the complementary subspace, equation (1) reduces
to multiplication by \(\alpha\). Thus the nonprojected part is retained
through \(T+R/\alpha\); it is not silently set to zero. For exact sample
cross-covariances, \(R\) is mathematically zero, but the implementation
retains the remainder to preserve the prepared numerical problem and
more general covariance inputs. For full bases the ordinary spectral
reconstruction is used.

The denominators in (3) change only when the augmentation parameter
changes. No Kronecker matrix of dimension
\((p_kp_\ell)\times(p_kp_\ell)\) is constructed, and no large inverse is
formed.

For example, computing \(Q_k^\top TQ_\ell\) can be associated in two
ways. The respective leading multiply counts are

\[
a_kp_kp_\ell+a_kp_\ell a_\ell,
\qquad
p_kp_\ell a_\ell+a_kp_ka_\ell.
\]

The code chooses the smaller count, and makes an analogous decision for
reconstruction. It uses double arithmetic for these cost estimates to
avoid integer overflow at large dimensions.

\hypertarget{algorithm-1-entrywise-proximal-and-dual-updates}{%
\section{Algorithm 1: entrywise proximal and dual
updates}\label{algorithm-1-entrywise-proximal-and-dual-updates}}

Once the coefficient solve is complete, the L11 auxiliary and scaled
dual updates are

\[
W=C+V,\qquad
Z^{+}=\operatorname{sign}(W)\bigl(|W|-\rho/\mu\bigr)_+,
\qquad V^{+}=W-Z^{+}.
\]

Each edge is processed independently. The compiled backend fuses
elementwise operations and residual accumulation, reducing temporary
arrays and R-level loop overhead. Its retained main state is still three
edge collections, or about \(3q\) doubles. The additional 40\% state
reduction described next applies to L21, not L11. The final reported L11
coefficients use the thresholded auxiliary blocks, with the established
numerical output cutoff.

\hypertarget{algorithm-2-lines-712-compressed-group-state}{%
\section{Algorithm 2, lines 7--12: compressed group
state}\label{algorithm-2-lines-712-compressed-group-state}}

\hypertarget{why-the-auxiliary-and-dual-arrays-can-share-storage}{%
\subsection{Why the auxiliary and dual arrays can share
storage}\label{why-the-auxiliary-and-dual-arrays-can-share-storage}}

For a row of a group proximal update, let

\[
W_k=M_k(C^{+})+V_k,\qquad
 a_{ki}^{+}=\left(1-\frac{\lambda_g}{\mu_g\|W_k[i,\cdot]\|_2}\right)_+.
\tag{5}
\]

A zero norm with positive penalty gives multiplier zero; zero penalty
gives multiplier one. Then

\[
G_k^{+}[i,\cdot]=a_{ki}^{+}W_k[i,\cdot],\qquad
V_k^{+}[i,\cdot]=(1-a_{ki}^{+})W_k[i,\cdot].
\]

Consequently, after a proximal/dual update it suffices to retain

\[
H_k=G_k+V_k,\qquad
G_k=\operatorname{diag}(a_k)H_k,\qquad
V_k=\operatorname{diag}(1-a_k)H_k.
\tag{6}
\]

These are exact identities for the updated state, not a low-rank
approximation. In particular,

\[
G_k-V_k=\operatorname{diag}(2a_k-1)H_k,
\tag{7}
\]

which supplies the next coefficient update directly.

The stored \texttt{Hk} and \texttt{Hl} are the two endpoint slices of
(6), both kept in the edge's \(p_k\times p_\ell\) orientation. The
multiplier for endpoint \(k\) scales rows; the multiplier for endpoint
\(\ell\) scales columns. This avoids repeatedly transposing the large
second-endpoint array.

\hypertarget{a-two-pass-subalgorithm}{%
\subsection{A two-pass subalgorithm}\label{a-two-pass-subalgorithm}}

The following is the implementation of the group update without
constructing \(M_k\), \(W_k\), \(G_k\), or \(V_k\) as full concatenated
matrices.

\begin{enumerate}
\def\labelenumi{\arabic{enumi}.}
\tightlist
\item
  Keep the old endpoint arrays \(H\) and row multipliers \(a\)
  unchanged. Initialize one squared-norm accumulator of length \(p_k\)
  for every view.
\item
  Traverse unordered edges. Form the target (2) using (7), and solve
  (1)--(4) for the new \(C_{k\ell}\).
\item
  On the same edge, form the transient endpoint blocks
\end{enumerate}

\[
W_{k\ell}^{(k)}=C_{k\ell}^{+}
 +\operatorname{diag}(1-a_k)H_{k\ell}^{(k)},
\qquad
W_{k\ell}^{(\ell)}=C_{k\ell}^{+}
 +H_{k\ell}^{(\ell)}\operatorname{diag}(1-a_\ell).
\tag{8}
\]

\begin{enumerate}
\def\labelenumi{\arabic{enumi}.}
\setcounter{enumi}{3}
\tightlist
\item
  Add the row sums of squared entries of \(W^{(k)}\) to view \(k\)'s
  accumulator, and the column sums for \(W^{(\ell)}\) to view \(\ell\)'s
  accumulator. Discard these temporary blocks. After all incident edges
  have contributed, take square roots and compute the new multipliers in
  (5).
\item
  Traverse the edges a second time. Recompute (8), recover the old and
  new auxiliary values transiently from (6), accumulate residuals, and
  overwrite the endpoint \(H\) arrays with their corresponding \(W\)
  blocks.
\item
  After every edge has been processed, replace \(a\) by \(a^{+}\). Check
  convergence and, when required, update the augmentation parameter.
\end{enumerate}

The second pass trades inexpensive rereading/recomputation for avoiding
persistent full auxiliary and dual matrices. Crucially, the shrinkage
norm includes \textbf{all} incident edges for a variable. Shrinking each
edge row separately would implement a different penalty.

\hypertarget{preserve-the-correct-stopping-rule}{%
\subsection{Preserve the correct stopping
rule}\label{preserve-the-correct-stopping-rule}}

With both endpoint slices in the edge orientation, the primal residual
is

\[
r_{\rm pri}^2=\sum_{k<\ell}
\left(\|C_{k\ell}-G_{k\ell}^{(k)}\|_F^2
+\|C_{k\ell}-G_{k\ell}^{(\ell)}\|_F^2\right).
\]

Writing \(\Delta G=G^{+}-G\), the dual residual is

\[
r_{\rm dual}^2=\mu_g^2\sum_{k<\ell}
\|\Delta G_{k\ell}^{(k)}+\Delta G_{k\ell}^{(\ell)}\|_F^2.
\tag{9}
\]

The endpoint changes must be added \textbf{before} squaring; replacing
(9) by the sum of their separate squared norms loses the cross term. The
code accumulates these scalars while visiting edges. The
absolute/relative tolerances remain

\[
\begin{aligned}
\epsilon_{\rm pri}
&=\sqrt{2q}\,\epsilon_{\rm abs}
+\epsilon_{\rm rel}\max\left\{
\left(2\sum_{k<\ell}\|C_{k\ell}\|_F^2\right)^{1/2},
\left(\sum_k\|G_k\|_F^2\right)^{1/2}\right\},\\
\epsilon_{\rm dual}
&=\sqrt q\,\epsilon_{\rm abs}
+\epsilon_{\rm rel}\mu_g
\left(\sum_{k<\ell}\|V_{k\ell}^{(k)}+V_{k\ell}^{(\ell)}\|_F^2\right)^{1/2}.
\end{aligned}
\]

Neither a small coefficient change nor a small step is substituted for
these residual tests.

\hypertarget{adaptive-penalties-and-old-warm-starts}{%
\subsection{Adaptive penalties and old warm
starts}\label{adaptive-penalties-and-old-warm-starts}}

If residual balancing changes \(\mu_g\) to \(f\mu_g\), the scaled dual
must become \(V/f\). For each row set

\[
b=a+\frac{1-a}{f},\qquad H\leftarrow bH,\qquad a\leftarrow\frac a b.
\tag{10}
\]

Indeed the new \(aH\) equals the old \(G\), and the new \((1-a)H\)
equals the old \(V/f\). This preserves the usual scaled-ADMM rescaling
rule without expanding the dual state; see S3.

An arbitrary legacy warm start need not initially satisfy (6). The code
accepts the old arrays, performs one ordinary group-proximal/dual
update, and then compresses the resulting state. It does not project an
arbitrary starting point into a different one. A coefficient-only start
uses the existing \(G=M(C)\), \(V=0\) convention, represented by
\(H=M(C)\) and \(a=1\).

\hypertarget{the-memory-reduction-and-its-limits}{%
\subsection{The memory reduction and its
limits}\label{the-memory-reduction-and-its-limits}}

Before compression, the principal L21 state contains \(C\) plus two
endpoint copies each of \(G\) and \(V\): \(5q\) doubles. After
compression it contains \(C\), two endpoint copies of \(H\), and \(p\)
row multipliers: \(3q+p\) doubles. Thus

\[
\text{fraction saved}=\frac{5q-(3q+p)}{5q}
=\frac25-\frac{p}{5q}\approx40\%.
\]

This counts the main state, not total process RAM. The large benchmark's
retained R state decreased from 96,802,984 to 58,104,864 bytes,
including object overhead. Covariance blocks, projected caches,
remainders, coefficient outputs, R/native conversion copies, workspaces,
and optional histories remain additional allocations. Group legacy
output requested with \texttt{compact\_state=FALSE} is reconstructed
only for retention after the optimized iterations.

For balanced views, \(q\) remains quadratic in \(p\). Exact entrywise or
row-group thresholding generally does not preserve a fixed low-rank
coefficient factorization. Therefore these changes do not justify a
claim that the whole estimator now uses \(O(pr)\) storage.

\hypertarget{algorithms-1-and-2-support-and-only-the-leading-r-eigenpairs}{%
\section{Algorithms 1 and 2: support and only the leading r
eigenpairs}\label{algorithms-1-and-2-support-and-only-the-leading-r-eigenpairs}}

\hypertarget{support-without-assembling-c}{%
\subsection{Support without assembling
C}\label{support-without-assembling-c}}

For each edge, add its squared row norms to view \(k\) and squared
column norms to view \(\ell\). This yields the full symmetric operator's
row norms without constructing that operator. The selected set is the
original

\[
\widehat S_\tau=\{i:\|\widehat C_{i\cdot}\|_2>\tau\}.
\]

Write \(s=|\widehat S_\tau|\). The existing numerical coefficient
cleanup and row threshold are retained.

\hypertarget{apply-metric-square-roots-instead-of-storing-them}{%
\subsection{Apply metric square roots instead of storing
them}\label{apply-metric-square-roots-instead-of-storing-them}}

Let \(\widetilde\Sigma_{0,\widehat S\widehat S}\) be the restricted
block covariance used for loading normalization. Its ridge is the
configured multiplier times the mean selected marginal variance;
eigenvalues are floored at \(10^{-10}\), as in the previous
implementation.

For a selected covariance block with basis \(Q\) and eigenvalues
\(d_j\), let \(b=\max(\xi,10^{-10})\) for an omitted nullspace, and let
\(\bar d_j=\max(d_j+\xi,10^{-10})\). Its action is

\[
\widetilde\Sigma^{\,\beta}x
=b^\beta x+Q\operatorname{diag}(\bar d_j^\beta-b^\beta)Q^\top x,
\qquad \beta\in\{1/2,-1/2\}.
\tag{11}
\]

For a full basis, the equivalent direct spectral action is used. The
code stores basis vectors and scalar weights, not dense square-root or
inverse-square-root matrices.

A full selected view reuses the prepared basis. For a partial support in
a wide view, it can factor the thin matrix
\(Q[\widehat S_k,\,\cdot]D^{1/2}\) rather than form a large selected
covariance. Otherwise it uses a selected covariance eigendecomposition.
These metric factorizations are distinct from the final eigenproblem;
the implementation does not promise that no full decomposition of any
small covariance is ever computed.

\hypertarget{the-final-operator-callback}{%
\subsection{The final operator
callback}\label{the-final-operator-callback}}

The required matrix is

\[
R_{\widehat S}
=\widetilde\Sigma_{0,\widehat S\widehat S}^{1/2}
 \widehat C_{\widehat S\widehat S}
 \widetilde\Sigma_{0,\widehat S\widehat S}^{1/2}.
\]

For any input vector \(v\), the callback performs three operations:
apply the right metric half-power using (11); multiply by selected
coefficient edges; apply the left metric half-power. On an edge, the
middle step accumulates

\[
y_k\mathrel{+}=C_{k\ell}[\widehat S_k,\widehat S_\ell]x_\ell,
\qquad
y_\ell\mathrel{+}=C_{k\ell}[\widehat S_k,\widehat S_\ell]^\top x_k.
\tag{12}
\]

The C++ routine reads those entries directly from the stored edge. It
does not allocate a new selected edge matrix on every callback. The
independent R backend uses zero-filled per-view vectors and BLAS
products with the stored edges.

The numerical request is conceptually

\begin{Shaded}
\begin{Highlighting}[]
\NormalTok{RSpectra}\SpecialCharTok{::}\FunctionTok{eigs\_sym}\NormalTok{(operator, }\AttributeTok{k =}\NormalTok{ r, }\AttributeTok{n =}\NormalTok{ s, }\AttributeTok{args =}\NormalTok{ context,}
                  \AttributeTok{which =} \StringTok{"LA"}\NormalTok{, }\AttributeTok{opts =} \FunctionTok{list}\NormalTok{(}\AttributeTok{tol =} \FloatTok{1e{-}12}\NormalTok{))}
\end{Highlighting}
\end{Shaded}

\texttt{"LA"} requests the largest \textbf{algebraic} eigenvalues. The
default \texttt{"LM"}, or a leading singular-value calculation, can
select large negative modes of this indefinite symmetric operator. They
are not interchangeable. RSpectra's documented function interface
supports exactly this matrix-vector representation; see S4.

The code asks for \(r\) output pairs but uses a larger internal Krylov
basis. Its first basis size is \(\min\{s,\max(4r+1,30)\}\); a retry uses
\(\min\{s,\max(8r+1,60)\}\). Approximate eigensolver storage is
\(O(sm+m^2+sr)\) for basis size \(m\), in addition to stored edges and
metric factors, rather than a newly assembled \(s^2\) operator. Runtime
depends on the number of matrix-vector products, spectral gaps, and
reorthogonalization; it is not universally \(O(sr)\).

Finite eigenpairs, the requested convergence count, residuals
\(\|Ru_j-\theta_ju_j\|_2\), and orthogonality are checked. An inadequate
result is retried with the larger Krylov space. Persistent failure
raises an error without an automatic dense EVD. The positivity
requirement is preserved, including rejection when the residual leaves
the positive-eigenvalue threshold unresolved. A two-dimensional problem
has an analytic leading-pair solution. Explicit dense mode, an
explicitly raised minimum partial-solve dimension, or requesting all
\(s\) pairs permits dense recovery; the reference backend also remains
dense for diagnostics.

Finally,

\[
\widehat L_{\widehat S}
=\widetilde\Sigma_{0,\widehat S\widehat S}^{-1/2}\widehat U_r,
\qquad \widehat L_{\widehat S^c}=0,
\qquad \widehat\lambda_j=1+\widehat\theta_j.
\]

\hypertarget{cross-validation-and-retained-outputs}{%
\section{Cross-validation and retained
outputs}\label{cross-validation-and-retained-outputs}}

For validation observations centered using \textbf{training} means,
partition \(L\) by views and compute \(Y_k=X_kL_k\). Then

\[
Q=\frac1{n_v}\sum_kY_k^\top Y_k=L^\top\widehat\Sigma_0L,
\qquad
A=\frac1{n_v}\left(\sum_kY_k\right)^\top
                   \left(\sum_kY_k\right)=L^\top\widehat\Sigma L.
\]

For \(n_v<p\), the validation object keeps the views rather than two
dense \(p\times p\) covariance matrices. All score matrices above are
only \(r\times r\). If \(Q=V\operatorname{diag}(d_j)V^\top\), the
existing ridge/floor gives \(\bar d_j\), and

\[
\operatorname{score}(L)=\sum_j\frac{v_j^\top Av_j}{\bar d_j}.
\]

This avoids an explicitly assembled inverse square root even in the
small score problem. CV minimizes its negative and still requires a
finite score on every fold.

Other memory controls bound the support-factor cache by both entry count
and estimated bytes (64 MiB by default), reuse fold preparation and warm
starts, and make histories/full operators/large retained fits optional.
The cache bound is not a total RAM limit. Independent CV workers have
private mutable solver states; supported fork workers can share
read-only pages until a write, while separate-session workers require
independent copies. More workers can therefore increase peak memory
substantially.

\hypertarget{storage-improvements-for-sgca-rgcca-sgcca-and-multicca}{%
\section{Storage improvements for SGCA, RGCCA, SGCCA and
MultiCCA}\label{storage-improvements-for-sgca-rgcca-sgcca-and-multicca}}

The following are \textbf{proposals based on the current wrappers and
upstream implementations}, not measured speedups or new numerical
backends in version 0.2.13. That version changes SGCA's time policy; it
does not silently replace the other estimators.

\begin{longtable}[]{@{}
  >{\raggedright\arraybackslash}p{(\columnwidth - 4\tabcolsep) * \real{0.3333}}
  >{\raggedright\arraybackslash}p{(\columnwidth - 4\tabcolsep) * \real{0.3333}}
  >{\raggedright\arraybackslash}p{(\columnwidth - 4\tabcolsep) * \real{0.3333}}@{}}
\toprule\noalign{}
\begin{minipage}[b]{\linewidth}\raggedright
Method
\end{minipage} & \begin{minipage}[b]{\linewidth}\raggedright
Main opportunity
\end{minipage} & \begin{minipage}[b]{\linewidth}\raggedright
What remains or must be preserved
\end{minipage} \\
\midrule\noalign{}
\endhead
\bottomrule\noalign{}
\endlastfoot
SGCA (Gao--Ma) & Data-matrix products in gradient descent; structured
metric operations in initialization & Dense lifted initializer and its
projection; original thresholding and backtracking \\
RGCCA & Remove wrapper Gram copies; reuse appropriate small metric
factors & Existing automatic primal/dual choice and parameter-dependent
deflation \\
SGCCA & Work through scores and feature-weight vectors; discard large
raw fit objects when unnecessary & Feature-space sparsity constraint,
scheme, sequential updates and deflation \\
MultiCCA & Data products and low-rank deflation corrections instead of
several dense Gram collections & PMA update order, threshold helper,
scaling and stopping convention \\
\end{longtable}

\hypertarget{sgca-gradient-descent-and-initialization-are-different-problems}{%
\subsection{SGCA: gradient descent and initialization are different
problems}\label{sgca-gradient-descent-and-initialization-are-different-problems}}

The current gradient code holds dense \(A=\widehat\Sigma\) and
\(B=\widehat\Sigma_0+\delta I\). Its scaled-iterate objective is

\[
f(W)=-\frac12\operatorname{tr}(W^\top AW)
+\frac14\|W^\top BW-\lambda I_r\|_F^2,
\qquad
\nabla f(W)=-AW+BW(W^\top BW-\lambda I_r).
\]

There is an exact data-space implementation. Partition \(W\) by views
and calculate

\[
(AW)_k=\frac1nX_k^\top\left(\sum_\ell X_\ell W_\ell\right),
\qquad
(BW)_k=\frac1nX_k^\top(X_kW_k)+\delta W_k.
\tag{13}
\]

No concatenated \(X\), full \(A\), or full \(B\) is needed for these
products. For small \(r\), their work is approximately \(O(npr)\) and
the extra working arrays can be \(O(pr+nr+r^2)\) beyond the data. The
row hard threshold, objective test, gradient-mapping stopping rule and
backtracking can remain unchanged. Gao--Ma also describe data-level
multiplication for their gradient step; see S5. This improvement matters
most when \(p\gg n\); a precomputed covariance can be preferable for
other shapes or many repeated products.

The current \textbf{initializer} is different: \texttt{Pi}, \texttt{H},
\texttt{Gamma}, and previous/temporary arrays have dimension
\(p\times p\), and \texttt{updateH()} performs a dense symmetric EVD.
Replacing gradient products alone does not remove that peak.

One exact restructuring is to retain \(B\) and \(B^{1/2}\) as view
blocks, so products such as \(B\Pi B\) are computed blockwise. Another
is to retain only loading factors and scalar diagnostics after
initialization, rather than every dense intermediate. Reordering the
candidate traversal to finish all \((k,\lambda)\) values for one
\(\rho\) before releasing that initializer could bound caches; candidate
identifiers and tie rules would need to be preserved. Some caching and
optional raw-fit retention already exist.

A detail in the bundled source matters: its projection clips shifted
eigenvalues as

\[
h_i=\min\{1,\max(d_i-\theta,0)\}.
\]

It uses the constraint \(\sum_i h_i\leq r\) and leaves \(\theta=0\) when
already feasible. The paper's Fantope is defined with trace
\textbf{equal} to \(r\). This pre-existing difference was preserved, not
changed as a performance edit. In either case, the projected matrix may
have more than \(r\) positive eigenvalues, so computing only \(r\)
eigenpairs is not generally an exact projection. A correct partial
projection would have to expand the spectral search until the omitted
spectrum is certified inactive; its rank cannot be fixed at \(r\) in
advance. A fixed rank factorization or a different initializer would
change the method and require separate statistical validation.

Also, the current extraction from \texttt{Pi} uses a leading
\textbf{SVD}. Replacing it with a partial SVD can preserve that rule,
subject to numerical checks. Replacing it with largest-algebraic
eigenvectors can change the result when a finite-iteration \texttt{Pi}
is indefinite. This is separate from EGCAR's final symmetric loading
EVD.

\hypertarget{rgcca-use-the-dual-formulation-already-available}{%
\subsection{RGCCA: use the dual formulation already
available}\label{rgcca-use-the-dual-formulation-already-available}}

The public RGCCA solver already selects a primal or dual formulation
according to block dimensions and reports \texttt{primal\_dual}; see S6.
For a wide block this moves regularized metric work to an \(n\times n\)
matrix. Therefore implementing another dual solver is not the first
priority.

For intuition, with \(M=\tau I+(1-\tau)X^\top X/n\) and \(0<\tau\leq1\),
Woodbury gives

\[
M^{-1}g=\tau^{-1}\left[g-cX^\top(I+cXX^\top)^{-1}Xg\right],
\qquad c=\frac{1-\tau}{n\tau}.
\]

A factorization can apply the smaller solve without forming an inverse.
The \(\tau=0\) case needs the intended singular-system/pseudoinverse
treatment. Reusing a fixed block's spectrum across a grid of positive
\(\tau\) values is possible, but later deflated blocks can depend on the
candidate and cannot be blindly reused.

The EGCAR wrapper currently builds \texttt{training\_gram\_blocks()} for
sign alignment even though it already has the centered training views.
Its existing fallback can align signs using the much smaller score
matrix \(Y=[X_1a_1,\ldots,X_Ka_K]\) and \(Y^\top Y/n\). Avoiding those
wrapper-level Gram arrays is a promising exact saving. Retaining only
\texttt{astar}, scores if needed, and scalar diagnostics can also reduce
saved output. \texttt{astar} must be preserved for mapping the original
centered data to all deflated components; substituting the raw deflated
weights \texttt{a} is not equivalent.

\hypertarget{sgcca-exploit-score-vectors-retain-feature-space-sparsity}{%
\subsection{SGCCA: exploit score vectors, retain feature-space
sparsity}\label{sgcca-exploit-score-vectors-retain-feature-space-sparsity}}

In the inspected RGCCA source, a sparse block uses a feature-weight
update and the specified constraint projection. The common inner loop
forms directions from block scores. For a direction
\(z_k\in\mathbb R^n\), its feature update uses \(X_k^\top z_k\), which
requires no \(p\times p\) cross-covariance. Thus storing data, scores,
and a few feature-weight vectors is sufficient for this part of the
calculation; see S7.

The same sign-alignment and output-retention improvements proposed for
RGCCA apply. However, an ordinary low-dimensional dual
reparameterization does not automatically preserve the featurewise L1
constraint: thresholded feature weights need not stay in the original
row space. Any custom backend must preserve the sparsity-bound scaling,
factorial/centroid/horst scheme, sequential update order, and component
deflation. We should benchmark the current upstream implementation
before rewriting it.

\hypertarget{multicca-avoid-repeated-gram-and-deflation-matrices}{%
\subsection{MultiCCA: avoid repeated Gram and deflation
matrices}\label{multicca-avoid-repeated-gram-and-deflation-matrices}}

The current \texttt{multicca\_common\_cv()} prepares normalized, raw and
centered Gram collections. Its Gram backend also forms deflated
cross-products for later components. Since the training views are
already centered, the centered correction is mathematically redundant,
apart from floating-point centering residuals. Storing transposed pair
blocks duplicates further data. Even choosing
\texttt{multicca\_backend="PMA"} currently passes through this shared
preparation, so changing that option alone does not eliminate wrapper
Gram storage.

An exact alternative applies cross-products through data:

\[
(X_i^\top X_j)w_j=X_i^\top(X_jw_j).
\]

For previously extracted weight matrices \(W_i,W_j\) and diagonal
deflation coefficients \(D_{ij}\), use

\[
\left(X_i^\top X_j-W_iD_{ij}W_j^\top\right)w_j
=X_i^\top(X_jw_j)-W_iD_{ij}(W_j^\top w_j).
\tag{14}
\]

This retains low-rank deflation factors instead of a new dense
cross-product. Scores \(X_jw_j\) can be updated as weights change. For a
fully connected graph, the nondeflated sums can be formed from a shared
sum of scores, subtracting the current block, while respecting the
sequential sweep. PMA's upstream \texttt{UpdateW} already expresses its
update through data and deflation factors; see S8.

For \(p\gg n\) this can exchange \(O(p^2)\) Gram storage for data plus
approximately \(O(pr)\) retained weights and small score arrays. It is a
memory/runtime tradeoff, not a universal speed improvement. The existing
Gram backend can win for modest \(p\) and many candidates. A useful
future interface would choose a backend using estimated bytes, with
numerical comparisons of loadings, criteria and component deflation. The
present PMA diagnostic iteration count is not by itself a convergence
certificate.

\hypertarget{a-common-prerequisite}{%
\subsection{A common prerequisite}\label{a-common-prerequisite}}

The shared \texttt{egcar\_cv\_data} object currently prepares EGCAR
training covariances for every fold. Even a fully data-based comparison
backend would still carry these if it used that same object unchanged.
End-to-end storage improvement therefore also calls for a lightweight
fold representation or lazy method-specific preparation. Shared fold
indices, training means, validation scoring, seeds, penalties, and tie
rules should remain identical. Reducing the worker count and avoiding
large retained fits are immediately available experiment controls;
changing grids or convergence tolerances is not an equivalent storage
optimization.

\hypertarget{sgcas-six-hour-policy-in-version-0.2.13}{%
\section{SGCA's six-hour policy in version
0.2.13}\label{sgcas-six-hour-policy-in-version-0.2.13}}

The new default is one \textbf{elapsed wall-clock budget of 21,600
seconds} for a complete SGCA CV grid and its full-sample refit. This is
not CPU time and not six hours for each candidate. A new budget starts
for each dataset/rank/replicate in the experiment runner. Shared
data/fold preparation happens first; SGCA-specific preparation,
initialization, candidate evaluation, fold work and final refit are
inside the budget.

\begin{Shaded}
\begin{Highlighting}[]
\NormalTok{bm }\OtherTok{\textless{}{-}} \FunctionTok{benchmark\_control}\NormalTok{(}
  \AttributeTok{sgca\_time\_limit =} \DecValTok{6} \SpecialCharTok{*} \DecValTok{60} \SpecialCharTok{*} \DecValTok{60}\NormalTok{,}
  \AttributeTok{sgca\_init\_max\_iter =} \ConstantTok{Inf}\NormalTok{,}
  \AttributeTok{sgca\_tgd\_max\_iter =} \ConstantTok{Inf}
\NormalTok{)}
\NormalTok{fit }\OtherTok{\textless{}{-}} \FunctionTok{sgca\_cv}\NormalTok{(shared, }\AttributeTok{rank =} \DecValTok{1}\NormalTok{, }\AttributeTok{benchmarks =}\NormalTok{ bm)}

\NormalTok{cfg }\OtherTok{\textless{}{-}} \FunctionTok{egcar\_experiment\_config}\NormalTok{(}\AttributeTok{sgca\_time\_limit =} \DecValTok{6} \SpecialCharTok{*} \DecValTok{60} \SpecialCharTok{*} \DecValTok{60}\NormalTok{)}
\end{Highlighting}
\end{Shaded}

Both SGCA iteration caps now default to \texttt{Inf}. They can still be
set explicitly for reproducibility or short tests. The TGD iteration
uses a \texttt{while} loop, so an infinite limit does not allocate an
infinite index sequence. Its existing stopping test still checks both
relative gradient mapping and relative iterate change; a tiny
backtracking step alone is not convergence.

A parent process monitors a \texttt{callr} R subprocess, polling at most
every 0.1 seconds during execution. On budget exhaustion it stops the
process tree, including fold workers. This can interrupt an initializer
blocked in compiled linear algebra, which a check only between R
iterations cannot do. Startup, serialization, scheduling and termination
introduce overhead, so the returned elapsed duration is measured rather
than assumed to be exactly 21,600 seconds. The subprocess holds its own
input copy: this timeout mechanism can increase peak memory and is not
one of the claimed EGCAR memory reductions. Setting the budget to
\texttt{Inf} runs in-process without supervision. See S9 for the process
API.

\begin{longtable}[]{@{}
  >{\raggedright\arraybackslash}p{(\columnwidth - 4\tabcolsep) * \real{0.3333}}
  >{\raggedright\arraybackslash}p{(\columnwidth - 4\tabcolsep) * \real{0.3333}}
  >{\raggedright\arraybackslash}p{(\columnwidth - 4\tabcolsep) * \real{0.3333}}@{}}
\toprule\noalign{}
\begin{minipage}[b]{\linewidth}\raggedright
Outcome
\end{minipage} & \begin{minipage}[b]{\linewidth}\raggedright
Reported status
\end{minipage} & \begin{minipage}[b]{\linewidth}\raggedright
Convergence and loading
\end{minipage} \\
\midrule\noalign{}
\endhead
\bottomrule\noalign{}
\endlastfoot
Procedure finishes and final stopping rules pass & \texttt{ok} &
\texttt{converged=TRUE}, loading returned \\
Explicit finite iteration cap is reached & \texttt{not\_converged} &
\texttt{converged=FALSE}; a valid approximate loading may be returned \\
Six-hour budget is exhausted & \texttt{time\_limit} &
\texttt{timed\_out=TRUE}, \texttt{converged=FALSE}, \texttt{L=NULL} \\
A numerical fit/CV error occurs & Existing error status, or propagated
error & Not relabeled as a timeout \\
\end{longtable}

The timeout message starts \textbf{``SGCA did not converge in real
time''} and states the budget. More precisely, the requested
CV-and-refit procedure did not finish in the allowed time. It does not
prove mathematical nonconvergence. An incomplete tuning run does not
choose the best candidate seen so far or score unfinished loadings as a
successful result. A watchdog timeout currently returns empty CV tables
and records whether it interrupted CV/startup or refitting; it does not
retain every interrupted candidate's state.

The installed package entry points implement this policy. The historical
script \texttt{run\_EGCAR\_local.R} in \texttt{inst/standalone/} is an
independent reference and retains its old controls. Use
\texttt{sgca\_cv()} or \texttt{run\_egcar\_experiments()} for the new
policy.

\hypertarget{evidence-reproduction-and-limits}{%
\section{Evidence, reproduction and
limits}\label{evidence-reproduction-and-limits}}

\hypertarget{previous-egcar-measurements}{%
\subsection{Previous EGCAR
measurements}\label{previous-egcar-measurements}}

The smaller test used \(n=50\), \((p_1,p_2,p_3)=(320,360,400)\), \(r=2\)
and 40 iterations. The larger used \(n=30\), \((800,900,1000)\), \(r=2\)
and 20 iterations. Both used direct penalty 0.02, prepared inputs and
the serial compiled backend.

\begin{longtable}[]{@{}
  >{\raggedleft\arraybackslash}p{(\columnwidth - 8\tabcolsep) * \real{0.2000}}
  >{\raggedright\arraybackslash}p{(\columnwidth - 8\tabcolsep) * \real{0.2000}}
  >{\raggedleft\arraybackslash}p{(\columnwidth - 8\tabcolsep) * \real{0.2000}}
  >{\raggedleft\arraybackslash}p{(\columnwidth - 8\tabcolsep) * \real{0.2000}}
  >{\raggedleft\arraybackslash}p{(\columnwidth - 8\tabcolsep) * \real{0.2000}}@{}}
\toprule\noalign{}
\begin{minipage}[b]{\linewidth}\raggedleft
Total p
\end{minipage} & \begin{minipage}[b]{\linewidth}\raggedright
Method
\end{minipage} & \begin{minipage}[b]{\linewidth}\raggedleft
Total seconds, 0.2.11
\end{minipage} & \begin{minipage}[b]{\linewidth}\raggedleft
Total seconds, 0.2.12
\end{minipage} & \begin{minipage}[b]{\linewidth}\raggedleft
Old/new ratio
\end{minipage} \\
\midrule\noalign{}
\endhead
\bottomrule\noalign{}
\endlastfoot
1080 & L11 & 1.430 & 1.400 & 1.02 \\
1080 & L21 & 1.456 & 1.333 & 1.09 \\
2700 & L11 & 4.100 & 2.197 & 1.87 \\
2700 & L21 & 5.577 & 2.980 & 1.87 \\
\end{longtable}

At \(p=2700\), L11 loading extraction decreased from 2.113 to 0.151
seconds and L21 extraction from 3.398 to 0.832 seconds. ADMM solve time
changed little. Both versions already used partial eigenpair extraction
at these dimensions; the key new difference was avoiding the dense
operator and its metric-root matrices, not merely changing a full EVD to
a partial one.

Across these benchmark cases, selected sets matched exactly. Maximum
coefficient, eigenvalue and eigenvector-projector differences were
approximately \(1.1\times10^{-17}\), \(1.6\times10^{-14}\) and
\(6.5\times10^{-14}\). Projectors allow for sign changes or basis
rotations within the recovered eigenspace.

Whole-process peak RSS for the large benchmark decreased from 750,192 to
422,368 KiB. That process included preparation and both penalty
families; these are not isolated per-fit allocation measurements. The
benchmark is too narrow to establish a general speed ratio, a comparison
against SGCA/RGCCA/SGCCA/MultiCCA, or a six-hour convergence rate.

\begin{Shaded}
\begin{Highlighting}[]
\ExtensionTok{Rscript}\NormalTok{ tools/benchmark\_0212.R OLD\_LIBRARY NEW\_LIBRARY results 3}
\ExtensionTok{Rscript}\NormalTok{ tools/benchmark\_0212.R OLD\_LIBRARY NEW\_LIBRARY results\_large 3 large}
\end{Highlighting}
\end{Shaded}

Version 0.2.12 passed 767 test expectations;
\texttt{R\ CMD\ check\ -\/-no-manual} reported zero errors and warnings
and three informational notes. One optional comparison suite was
skipped. The tests included independent dense comparisons, tall/wide and
degenerate inputs, old and compressed warm starts, adaptive penalties,
operator products, metric actions, algebraic eigenvalue order and
partial-solver failure. The runtime was R 4.3.3 on Ubuntu 24.04, x86-64,
with the system BLAS/LAPACK.

\hypertarget{new-timeout-validation}{%
\subsection{New timeout validation}\label{new-timeout-validation}}

Version 0.2.13 adds tests for default/invalid budgets,
unlimited-iteration numerical stopping, successful subprocess return,
solver errors, short actual elapsed-time termination, process-tree
cleanup, supervised/in-process agreement, serial/parallel SGCA
agreement, the reference initializer, and public timeout status
propagation. Short budgets test the same watchdog; \textbf{no six-hour
wait or full six-hour convergence experiment was performed}. The
complete updated suite passed \textbf{815 expectations}, with zero
failures or test warnings and one optional suite skipped.
\texttt{R\ CMD\ check\ -\/-no-manual} reported zero errors, zero
warnings and four informational notes: optional suggestions unavailable,
the C++14 specification, binary size, and installation of the requested
LaTeX source. The accompanying \texttt{VALIDATION\_0213.md} records the
details.

The proposed storage backends for the four comparison methods have not
been implemented or benchmarked in this release. They should be
validated against the existing code before using their timings in a
methods comparison.

\hypertarget{code-map}{%
\section{Code map}\label{code-map}}

\begin{longtable}[]{@{}
  >{\raggedright\arraybackslash}p{(\columnwidth - 2\tabcolsep) * \real{0.5000}}
  >{\raggedright\arraybackslash}p{(\columnwidth - 2\tabcolsep) * \real{0.5000}}@{}}
\toprule\noalign{}
\begin{minipage}[b]{\linewidth}\raggedright
Responsibility
\end{minipage} & \begin{minipage}[b]{\linewidth}\raggedright
Source location
\end{minipage} \\
\midrule\noalign{}
\endhead
\bottomrule\noalign{}
\endlastfoot
Covariance preparation, exact spectral solve support, metric factors and
callback & \texttt{R/helpers.r} \\
L11 dispatch and independent R ADMM implementation &
\texttt{R/reduced\_rank\_regression.R} \\
Group output, compatibility and retained state &
\texttt{R/group\_reduced\_rank\_regression.R} \\
Fused compiled ADMM and selected-edge multiplication &
\texttt{src/egcar\_native.cpp} \\
Common held-out score & \texttt{R/metrics.R} \\
SGCA initialization, cached algebra and penalized TGD &
\texttt{R/sgca\_initializer.R}, \texttt{R/alt\_SGCA.R} \\
Six-hour supervisor and timeout result & \texttt{R/sgca\_budget.R} \\
Public controls and experiment defaults & \texttt{R/controls.R},
\texttt{R/experiment\_config.R} \\
RGCCA/SGCCA and MultiCCA adapters & \texttt{R/alt\_RGCCA.R},
\texttt{R/alt\_SGCCA.R}, \texttt{R/alt\_MultiCCA\_CrossValidation.R} \\
Measured 0.2.12 benchmarks and raw evidence &
\texttt{tools/benchmark\_0212.R},
\texttt{inst/validation/performance\_0212\_*} \\
\end{longtable}

\hypertarget{sources}{%
\section{Sources}\label{sources}}

The mathematical reductions above are derived from the supplied
algorithms and checked against this package's source. External
implementation/documentation sources were inspected on 24 September
2026.

S1. ccar3 0.1.2, CRAN manual:
\url{https://cran.r-project.org/web/packages/ccar3/ccar3.pdf}.

S2. ccar3, source \texttt{R/ecca.r}:
\url{https://github.com/cran/ccar3/blob/master/R/ecca.r}.

S3. Boyd et al., \emph{Distributed Optimization and Statistical Learning
via the Alternating Direction Method of Multipliers}:
\url{https://web.stanford.edu/~boyd/papers/pdf/admm_distr_stats.pdf}.

S4. RSpectra manual, function interface and eigenvalue selection:
\url{https://cran.r-project.org/web/packages/RSpectra/RSpectra.pdf}.

S5. Gao and Ma (2023), \emph{Sparse GCA and Thresholded Gradient
Descent}, JMLR 24(135), 1--61:
\url{https://jmlr.org/papers/v24/21-0745.html}. Bundled initializer
provenance:
\url{https://github.com/TannisthaM/SGCA/blob/main/R/gao_cv_functions.R}.

S6. RGCCA public interface, output fields and primal/dual description:
\url{https://rgcca-factory.github.io/RGCCA/reference/rgcca.html}.

S7. RGCCA source, \texttt{block.R}, \texttt{block\_init.R},
\texttt{block\_update.R}, \texttt{rgcca\_inner\_loop.R}:
\url{https://github.com/rgcca-factory/RGCCA/tree/main/R}.

S8. PMA source, \texttt{R/MultiCCA.R}, particularly \texttt{UpdateW}:
\url{https://github.com/cran/PMA/blob/master/R/MultiCCA.R}.

S9. callr background R sessions:
\url{https://callr.r-lib.org/reference/r_bg.html}; processx process-tree
termination: \url{https://processx.r-lib.org/reference/process.html}.

\end{document}
